\clearpage\pagestyle{otMain}\setlength{\otBlocEspacement}{2.5mm}\fontsize{11}{15}\selectfont\corrigesChapitre{15}{Programmes de construction}
\begin{blocLivre}{cor-15-1}\begin{corrige}{1}{Les bons mots}\demarche{Associer chaque notation à son objet.}
Le segment $\res{[RS]}$, la droite $\res{(TU)}$ et la longueur $\res{RS}$, sans crochets : trois écritures pour trois objets.
\label{fig-cor-15-1}\figureGeo[corrige]{geo-1}
\end{corrige}\end{blocLivre}
\begin{blocLivre}{cor-15-2}\begin{corrige}{2}{Tracer une perpendiculaire}\demarche{Garder un côté de l’équerre sur la droite.}
L’équerre glisse le long de (d) jusqu’à M ; le tracé vert coupe (d) à angle droit codé.
\label{fig-cor-15-2}\figureGeo[corrige]{geo-2}
\end{corrige}\end{blocLivre}
\begin{blocLivre}{cor-15-3}\begin{corrige}{3}{Tracer une parallèle}\demarche{Bloquer la règle, faire glisser l’équerre.}
L’équerre s’appuie sur (d), la règle tient l’équerre, on glisse jusqu’à N : la parallèle verte garde le même écart.
\label{fig-cor-15-3}\figureGeo[corrige]{geo-3}
\end{corrige}\end{blocLivre}
\begin{blocLivre}{cor-15-4}\begin{corrige}{4}{Rayon et diamètre}\demarche{Le diamètre vaut deux rayons.}
Diamètre : $3{,}5\times2=\res7$ cm. Rayon : $9:2=\res{4{,}5}$ cm.
\label{fig-cor-15-4}\figureGeo[corrige]{geo-4}
\end{corrige}\end{blocLivre}
\begin{blocLivre}{cor-15-5}\begin{corrige}{5}{Exécuter un programme}\demarche{Suivre les étapes dans l’ordre écrit.}
Le rectangle se construit à l’équerre ; M est à $2{,}5$ cm de A ; le cercle de centre M passant par A a un rayon de $\res{2{,}5}$ cm et passe aussi par B. La construction figure en vert.
\label{fig-cor-15-5}\figureGeo[corrige]{geo-5}
\end{corrige}\end{blocLivre}
\begin{blocLivre}{cor-15-6}\begin{corrige}{6}{Le triangle du compas}\demarche{Relire la définition du cercle pour conclure.}
Les deux arcs de même rayon donnent GE $=$ GF $=4$ cm : le triangle est $\res{\text{isocèle}}$ en G. La construction verte garde les arcs visibles.
\label{fig-cor-15-6}\figureGeo[corrige]{geo-6}
\end{corrige}\end{blocLivre}
\begin{blocLivre}{cor-15-7}\begin{corrige}{7}{Le programme en désordre}\demarche{Chercher ce dont chaque phrase a besoin.}
L’ordre est $\res{\text{b, c, d, a}}$ : impossible de tracer le cercle de centre O avant d’avoir obtenu O. Chaque étape utilise les objets déjà construits.
\label{fig-cor-15-7}\figureGeo[corrige]{geo-7}
\end{corrige}\end{blocLivre}
\begin{blocLivre}{cor-15-8}\begin{corrige}{8}{L’étape impossible}\demarche{Confronter chaque phrase à la définition du cercle.}
La phrase 3 : un point du cercle est à exactement $\res{3\ \mathrm{cm}}$ de O, or P est à 5 cm, donc P est $\res{\text{à l’extérieur}}$ du cercle. On corrige : « P est à l’extérieur du cercle. »
\label{fig-cor-15-8}\figureGeo[corrige]{geo-8}
\end{corrige}\end{blocLivre}
\begin{blocLivre}{cor-15-9}\begin{corrige}{9}{Écrire le programme}\demarche{Nommer les points et donner les mesures.}
Un programme correct : trace [RS] de 3 cm ; trace la perpendiculaire à (RS) en R et place U à 3 cm de R ; trace la perpendiculaire à (RS) en S et place T à 3 cm de S ; joins U à T ; trace le segment [RT]. D’autres ordres corrects existent : l’essentiel est de ne rien utiliser avant de l’avoir construit.
\label{fig-cor-15-9}\figureGeo[corrige]{geo-9}
\end{corrige}\end{blocLivre}
\begin{blocLivre}{cor-15-10}\begin{corrige}{10}{Le cercle bien défini}\demarche{Définir le rayon par les deux points donnés.}
Programme : pointe du compas en O, ouverture jusqu’à A ; tracer le cercle. Le rayon reporté est la longueur $\res{OA}$ : le cercle passe alors exactement par A.
\label{fig-cor-15-10}\figureGeo[corrige]{geo-10}
\end{corrige}\end{blocLivre}
\begin{blocLivre}{cor-15-11}\begin{corrige}{11}{La figure au téléphone}\demarche{Chercher deux figures différentes qui respectent la consigne.}
$\res{\text{Max a raison}}$ : tous les losanges de côté 4 cm ne sont pas superposables, leur ouverture d’angle peut changer. Il faut ajouter une donnée, par exemple la longueur d’une diagonale ou un angle à décalquer.
\label{fig-cor-15-11}\figureGeo[corrige]{geo-11}
\end{corrige}\end{blocLivre}
\begin{blocLivre}{cor-15-12}\begin{corrige}{12}{Les cercles concentriques}\demarche{Soustraire les rayons voisins.}
Les trois cercles partagent le même centre : des cercles concentriques. Chaque anneau a une largeur de $2-1=\res1$ cm, puis $3-2=\res1$ cm : des anneaux réguliers.
\label{fig-cor-15-12}\figureGeo[corrige]{geo-12}
\end{corrige}\end{blocLivre}
\begin{blocLivre}{cor-15-13}\begin{corrige}{13}{La cible du stand}\demarche{Garder le même centre à chaque étape.}
Programme : place un point O ; trace le cercle de centre O et de rayon 2 cm ; puis de rayon 4 cm ; puis de rayon 6 cm. Le grand diamètre vaut $6\times2=\res{12}$ cm.
\label{fig-cor-15-13}\figureGeo[corrige]{geo-13}
\end{corrige}\end{blocLivre}
\begin{blocLivre}{cor-15-14}\begin{corrige}{14}{La maison à construire}\demarche{Construire le carré avant le toit.}
Le carré coiffé du triangle isocèle DTC dessine $\res{\text{une maison}}$ : le toit a son sommet T au-dessus du milieu M, donc DT $=$ CT. La construction figure en vert.
\label{fig-cor-15-14}\figureGeo[corrige]{geo-14}
\end{corrige}\end{blocLivre}
\begin{blocLivre}{cor-15-15}\begin{corrige}{15}{L’ordre qui compte}\demarche{Chercher une instruction qui dépend d’une autre.}
$\res{\text{Faux}}$ : « place le milieu M de [AB] » doit suivre le tracé de [AB] ; échangées, les instructions deviennent inexécutables. Certaines étapes indépendantes peuvent s’échanger, mais pas celles qui utilisent un objet encore absent.
\label{fig-cor-15-15}\figureGeo[corrige]{geo-15}
\end{corrige}\end{blocLivre}
